\chapter{Verification}
\label{ch:ch5_verification}

% [DRAFT v1] Source material: rtl/tb/*, rtl/run_tb.py, analysis/ctr_drbg_ref.py,
% analysis/verifica_top.py, firmware/verificador/, docs/mapa_registros.md.
% FPGA bring-up section pending hardware.

\section{Strategy: vectors, simulation, independent verifier}
\label{sec:verif_strategy}

A cryptographic engine that is verified only against itself proves
nothing: an error in the S-box would be reproduced identically in the
model used to check it. The verification of this work therefore rests on
three legs that do not share an implementation.

\begin{enumerate}
\item \textbf{Official vectors.} Every primitive is checked against the
test vectors published with its standard: FIPS~197 and SP~800-38A for
AES, RFC~4493 and the NIST example file for CMAC, and the CAVP response
files for CTR\_DRBG (with and without derivation function, with and
without reseed). These are known-answer tests and they settle
correctness of the algorithm.
\item \textbf{Independent reference models.} A bit-exact model of the
generator was written in Python with its own AES, derived from the
standard without looking at the VHDL. It is validated against the CAVP
vectors first, and only then used to generate expected values for the
fixed input layout of the hardware and to recompute the tags of the
system-level bench. Two implementations that agree are worth more than
one that passes.
\item \textbf{An independent verifier in the loop.} On hardware the
ESP32 recomputes with mbedtls, and with its own AES accelerator, every
result the engine produces, so that the final check is performed by a
third, unrelated implementation over a real bus.
\end{enumerate}

Two working rules apply throughout, both inherited from the author's
previous experimental work: the output of a compilation, simulation or
measurement is never filtered before being read in full --- a filter
hides exactly the warning that mattered --- and failures are recorded
whole, never discarded. Twice in this work a filter on the simulator's
log hid a result, and both times the fix was to remove the filter.

\section{Reference models in Python}
\label{sec:refmodels}

\texttt{ctr\_drbg\_ref.py} implements AES-128 encryption, the block-cipher
derivation function, CBC-MAC and CTR\_DRBG with and without derivation
function, in pure Python and from the text of SP~800-90A. The AES asserts
the two FIPS~197 vectors at import; the DRBG is exercised against the CAVP
response files following the CAVS procedure (instantiate, optional reseed,
two generate calls, compare the second): \textbf{960 AES-128 cases, all
passing}. One defect was found in the file parser rather than in the
model --- a section header must close the pending test case before the
parameters change, otherwise the last case of every parameter block is
evaluated with the next block's parameters --- and it is worth recording
because its symptom, a single failure at count~14 of each block, looked
exactly like an algorithmic error.

A second script, \texttt{gen\_drbg\_vectores.py}, drives the validated
model through the sequence the hardware supports (instantiate with 32
bytes of entropy and 16 of nonce, generate twice, reseed, generate) and
emits a VHDL package of expected values; if the input layout of the
hardware ever changes, that script is the only place to change it. A
third, \texttt{verifica\_top.py}, recomputes AES-CMAC from the same AES and
validates itself against RFC~4493 before checking anything.

\section{Testbenches and results}
\label{sec:testbenches}

All VHDL is simulated with GHDL~6 under a single runner that analyses the
sources in dependency order and executes every bench, printing the
simulator's output unfiltered. Table~\ref{tab:tb} summarizes the seven
benches and their state at the time of writing.

\begin{table}[htbp]
\centering
\caption{Testbenches, what they check and their result (GHDL 6.0, VHDL-2008).}
\label{tab:tb}
\begin{tabular}{@{}llc@{}}
\toprule
Bench & Checks & Result \\
\midrule
Ring oscillator & start/stop, $f = 1/(2N t_{pd})$, zero floor, $\sigma_T = \sqrt{2N}\,\sigma_{\mathrm{stage}}$ & 4/4 \\
I2C slave & addressing, foreign address, multi-byte write, repeated-start read & 14/14 \\
AES-128 core & FIPS~197 B, C.1; SP~800-38A F.1.1 (3 blocks); 1000-block chain & 6/6 \\
AES-CMAC & RFC~4493 examples 1--4; repeatability & 5/5 \\
CTR\_DRBG & instantiate, 2$\times$generate, reseed, counter, re-instantiate, refuse unseeded & 13/13 \\
Health tests & healthy source, stuck, 70\,\% biased, stuck-at-complement, sticky/clear & 10/10 \\
Whole engine & ESP32 sequence over I2C: id, seed, key, 2 challenges, 2nd key, capture, key hidden & 16/16 \\
\bottomrule
\end{tabular}
\end{table}

Some of what the benches found deserves to be stated, because a
verification chapter that reports only passes has not verified much.

\paragraph{Ring model.} The LUT primitives used to simulate the ring
were first modeled with a strict unknown-propagation rule: any unknown
input gave an unknown output. At the start of a simulation every signal
is unknown, so the ring never left that state and never started. A real
NAND with its enable input at zero outputs one regardless of the other
input; the model now evaluates every combination compatible with the
known inputs and resolves the output when they all agree. Without this
the entire system-level bench would have been impossible.

\paragraph{I2C slave.} Two genuine design errors: the register pointer
did not advance on multi-byte writes, so consecutive bytes overwrote the
same address; and on reads the most significant bit was driven one clock
edge too late, shifting every byte. The first was fixed by incrementing
the pointer one state after the write strobe, so the address is still
valid during the strobe; the second by loading the shift register and
driving its first bit on the same falling edge that releases the
acknowledge. A third defect was in the bench itself: initializing the
register memory from the stimulus process created a second driver on the
same signal and every conflicting bit resolved to unknown.

\paragraph{Health tests.} The ``healthy'' source of the first bench, a
shift register seeded with a single one, started with a run of 27 zeros
and a first window at 76\,\%; the detector fired. The detector was right.

\paragraph{Key leak through the random register.} The most serious
finding did not come from a bench but from reading the design again
before handing it over. The command that generated the key wrote the
full 256-bit output of the DRBG into the random-output register, which is
readable without the test pin; the first 128 of those bits \emph{were} the
key. The system bench even asserted that equality as a ``check''. The fix
separates the two: key generation leaves nothing readable, and a new
command fills the random register from an independent generation, after
the DRBG has updated its state. The bench now asserts the opposite, that
the register contains the key in neither half, and the Python
cross-check does the same.

\paragraph{Whole engine.} With the rings at their physical speed the
system bench, which drives almost a millisecond of I2C traffic, did not
finish in twenty minutes of CPU: every ring transition costs two random
draws in the jitter model. In that bench alone the stage delay is set ten
times larger, with the same relative jitter, and this is stated in the
bench. The sixteen checks pass, and the key, the two challenges and their
tags are written to a file that \texttt{verifica\_top.py} checks: both
tags equal the CMAC recomputed by the independent Python AES, the
readable random-output register contains the key in neither of its
halves, and two generations give different keys. This is the first point at which a key born from
ring-oscillator noise, conditioned, generated and used to authenticate,
is confirmed by an implementation that shares nothing with the VHDL.

\section{ESP32 verifier firmware}
\label{sec:firmware}

The firmware is an ESP-IDF~5 application for the classic ESP32 that
talks to the engine through the register map of Appendix~\ref{app:regs}
and offers a console over the UART. Its commands map onto the
verification plan: \texttt{id} reads identification and status;
\texttt{sembrar} and \texttt{generar} issue the long commands and poll
the busy flag; in test mode \texttt{generar} also reads back the key
(provisioning) and checks that it equals the first sixteen bytes of the
generated block; \texttt{reto N} performs $N$ challenge-response rounds
with nonces from \texttt{esp\_random()}, recomputing each tag with
\texttt{mbedtls\_cipher\_cmac} and counting matches and mismatches, the
mismatches being printed in full; \texttt{crudo N} dumps $N$ chunks of
raw bits and \texttt{esp N} dumps blocks of the ESP32's own generator, both
in the framed, CRC32-protected format used in the author's previous work
so that the same statistical pipeline processes the two sources. CMAC is
not enabled by default in the ESP-IDF build of mbedtls and is switched on
in the project defaults; the console runs at 115\,200~baud, the only rate
that proved reliable with the board's USB bridge in earlier campaigns,
and the log level is set to errors only so that log lines cannot corrupt
the frames. The firmware builds with ESP-IDF~5.3.2 into a 255~kB image;
running it awaits the hardware.

\section{FPGA bring-up}
\label{sec:bringup}

% [PENDING HARDWARE] Vivado project, constraints (DONT_TOUCH, pblocks,
% ALLOW_COMBINATORIAL_LOOPS, set_disable_timing), first bitstream: do the
% rings oscillate and at what frequency; then the ESP32 in the loop.
Pending hardware. The constraints the rings need are known from the
literature and from the vendor documentation --- \texttt{DONT\_TOUCH} on
the cells and nets of each ring (unlike \texttt{KEEP}, it survives
implementation), \texttt{ALLOW\_COMBINATORIAL\_LOOPS} on one net of each
loop, \texttt{set\_disable\_timing} to break the loop for static timing
analysis, and one placement block per ring so that identical rings are
neither merged nor placed side by side --- and the first bitstream has a
single purpose: to read the frequency counters and confirm that every
ring oscillates. Nothing about entropy is claimed before that.
